\documentclass[10pt]{article}
\usepackage[margin=0.72in]{geometry}
\usepackage[T1]{fontenc}
\usepackage[utf8]{inputenc}
\usepackage{array}
\usepackage{booktabs}
\usepackage{caption}
\usepackage{graphicx}
\usepackage{hyperref}
\usepackage{longtable}
\usepackage{ragged2e}
\usepackage{xcolor}
\usepackage{url}

\hypersetup{
  colorlinks=true,
  linkcolor=blue,
  urlcolor=blue,
  citecolor=blue
}

\newcommand{\statuskeep}{\textbf{Adopt}}
\newcommand{\statusmodify}{\textbf{Adopt with controls}}
\newcommand{\statusskip}{\textbf{Do not adopt blindly}}
\newcommand{\statusalready}{\textbf{Already better in 3DTI target}}
\newcolumntype{P}[1]{>{\RaggedRight\arraybackslash}p{#1}}

\title{DynaSpace Smartphone PPS Looming Stimulus\\Spectral and Spatial Feature Audit}
\author{PPS-kit derived report}
\date{2026-06-29}

\begin{document}
\maketitle

\section*{Executive Decision}

The smartphone DynaSpace stimulus should not be copied as an opaque WAV into PPS-kit. It should be decomposed into a parameterized dry burst train plus a 3DTI/SOFA trajectory renderer. The useful smartphone-study features are: a broadband Gaussian white-noise burst train, a strong but not extreme rising distance/intensity envelope, the 6.4 m to 0.2 m constant-bearing trajectory, log-spaced tactile anchors, and room/externalization cues visible in the low interaural coherence of the baked WAV. Those are worth adopting as explicit parameters.

The features that PPS-kit/3DTI should handle instead of imitating by raw spectral hacks are HRTF spectral filtering, ILD, ITD, near-field corrections, future trajectory angles, and any controlled reverberation/direct-to-reverberant behavior. The initial PPS-kit proxy audit found a usable but imperfect placeholder: it was continuous rather than burst-train-like, much higher in RMS level, less bright, highly interaurally coherent, and had an opposite-signed ILD/ITD convention relative to the DynaSpace raw files. The implementation decision documented here is therefore a best-of-both-worlds protocol: adopt the DynaSpace dry-source burst train as a source profile, but keep spatial texture and future trajectory flexibility inside the 3DTI/SOFA rendering path.

\section*{Scope and Inputs}

This audit isolates the raw spectral, temporal, level, and binaural features in the DynaSpace Android audio files used by Roussel et al.'s smartphone PPS validation study \cite{roussel2025}. The original raw WAV files were measured from a private local archival clone of the public DynaSpace source repository \cite{dynaspace_repo}. The raw files are not redistributed here. PPS-kit artifacts in this folder are derived metrics, plots, TeX, and PDF only.

The compared audio files are:
\begin{itemize}
  \item DynaSpace raw looming: \texttt{bursttrainlooming.wav}, SHA256 prefix \texttt{610c9a9c2232}. Full hash is stored in \path{dynaspace_audio_feature_metrics.json}.
  \item DynaSpace raw fixed: \texttt{bursttrainfixe.wav}, SHA256 prefix \texttt{56f8afff1b09}. Full hash is stored in \path{dynaspace_audio_feature_metrics.json}.
  \item PPS-kit proxy looming: \path{assets/preloads/roussel_2025_dynaspace_mobile_pps/02_looming_stimuli/looming_DynaSpace_proxy_looming_burst_train.wav}.
  \item PPS-kit proxy fixed: \path{assets/preloads/roussel_2025_dynaspace_mobile_pps/02_looming_stimuli/looming_DynaSpace_proxy_fixed_640_cm_burst_train.wav}.
\end{itemize}

\section*{Search and Consensus Method}

Search dates: 2026-06-28 to 2026-06-29. I used web searches across the smartphone PPS paper, DynaSpace app, Hobeika burst-train lineage, Canzoneri PPS looming task, auditory looming intensity literature, auditory localization/distance reviews, and 3DTI documentation. I then manually accessed Consensus through the user's browser and ran targeted searches that attempted to prove or disprove the assumptions behind each candidate feature. Those Consensus searches are exported in this report as a separate LaTeX section with re-openable search URLs.

\begin{longtable}{P{0.23\linewidth} P{0.34\linewidth} P{0.35\linewidth}}
\caption{Search iterations and evidence role.}\\
\toprule
Feature class & Search iteration & Evidence used in this report\\
\midrule
\endfirsthead
\toprule
Feature class & Search iteration & Evidence used in this report\\
\midrule
\endhead
Smartphone PPS parameters & Queries for Roussel, DynaSpace, 640 cm, 20 cm, 210 cm/s, LISTEN, and tactile anchors. & Roussel et al. report a DynaSpace Android app, a 3250 ms auditory stimulus, a fixed 640 cm source, a looming source from 640 cm to 20 cm at ear level and 210 cm/s, azimuth 45 degrees in left hemispace, and tactile delays T1--T6 at 105, 1625, 2385, 2765, 2955, and 3050 ms \cite{roussel2025}.\\
Burst-train source & Queries for Hobeika, Gaussian white noise bursts, 30 ms, 10 ms rise/fall, 65 ms intervals. & The lineage stimulus is a Gaussian white-noise burst train with short bursts and ramps; 30 ms bursts plus 65 ms interval predicts a 95 ms onset spacing, matching the 33 measured DynaSpace peaks \cite{hobeika2020}.\\
Basic PPS effect & Queries for dynamic sounds, boundaries of peripersonal space, approaching auditory stimuli. & Canzoneri et al. show that dynamic sounds can reveal spatially dependent audio-tactile facilitation and PPS boundaries, with approaching sounds producing the stronger tactile RT pattern \cite{canzoneri2012}.\\
Rising level and looming warning value & Queries for auditory looming bias, rising sound intensity, warning cue, amygdala, phasic alertness. & Rising intensity is supported as a warning/alerting cue and is central to auditory looming salience \cite{bach2008,neuhoff2016}.\\
Localization cue inventory & Queries for auditory localization ITD, ILD, HRTF, bandwidth, reverberation, direct-to-reverberant ratio. & Reviews identify ITD, ILD, HRTF spectral filtering, bandwidth, intensity, reverberation, and direct-to-reverberant ratio as spatial/distance cues \cite{carlini2024,kolarik2016}.\\
3DTI capability & Queries for 3D Tune-In Toolkit, SOFA, HRTF, distance attenuation, near-field, BRIR, reverberation. & 3DTI supports SOFA/HRIR convolution, ITD customization, distance attenuation, near-field corrections, movement smoothing, and BRIR/Ambisonic reverberation in the full toolkit \cite{cuevas2019}.\\
Local PPS-kit implementation & Repository search for FABIAN, SOFA, gain law, nearest HRIR, reverb, and generated proxy assets. & PPS-kit stores FABIAN as the standard redistributable HRTF resource, records LISTEN as excluded pending license audit, and the current Python reference renderer uses nearest HRIR lookup, 1/r gain when enabled, peak normalization unless loudness policy is enabled, and reverb disabled.\\
Consensus targeted checks & Browser-based Consensus searches for looming salience, burst/onset localization, HRTF/ILD/ITD, reverberation/DRR/IACC, near-field cues, PPS tactile RT effects, individualized HRTF value, and whether intensity alone is sufficient. & The Consensus pass supported adopting a burst-train source and keeping spatialization in 3DTI/SOFA, while warning that rising intensity and level imbalance can masquerade as spatial PPS effects unless matched-level controls are added.\\
\bottomrule
\end{longtable}

\section*{Source Summary}

\begin{longtable}{P{0.22\linewidth} P{0.34\linewidth} P{0.36\linewidth}}
\caption{Source-by-source role in the final decision.}\\
\toprule
Source group & What it contributed & What it did not justify\\
\midrule
\endfirsthead
\toprule
Source group & What it contributed & What it did not justify\\
\midrule
\endhead
Roussel et al. smartphone PPS paper and DynaSpace app & Established the runnable smartphone validation task: 3.25 s auditory stimulus inside a 3.85 s file, 6.4 m to 0.2 m left-hemispace trajectory at about 2.1 m/s, fixed 6.4 m comparator, six tactile anchors, no-vibration catches, Android playback, and app-level timing parameters. & It does not prove that the baked Android WAV is the best long-term stimulus representation for PPS-kit, because it fixes one trajectory and one non-individual HRTF/room render.\\
Raw DynaSpace WAV audit & Established the actual stimulus waveform properties: 33 bursts, about 95 ms inter-onset interval, clean headroom, high crest factor, bright broadband spectrum, rising active level for looming, low IACC consistent with non-anechoic spatial rendering, and a large fixed-vs-looming level difference. & It does not isolate which feature caused behavioral facilitation; the raw file combines source transients, intensity growth, HRTF, reverb, and level asymmetry.\\
Hobeika/Canzoneri/Serino PPS lineage & Justified dynamic approaching sounds and repeated burst-train auditory sources as established PPS-task ingredients, not arbitrary stimulus choices. & It does not remove expectancy and timing confounds; matched timing and level controls remain needed.\\
Auditory looming literature & Supported rising intensity as an intrinsic warning/salience cue and explained why looming sounds often speed detection or orienting. & It does not establish rising intensity as a pure spatial cue; intensity can drive effects without robust externalized distance perception.\\
Auditory localization/distance literature & Confirmed that broadband content, HRTF spectral filtering, ITD, ILD, near-field behavior, reverberation, and direct-to-reverberant ratio are the relevant spatial cue classes. & It does not support baking these cues into one static WAV when a parameterized renderer can represent them more transparently.\\
3DTI/SOFA documentation and PPS-kit code & Confirmed that the toolkit can keep the source waveform separate from spatial rendering and can expose trajectory/HRTF/level metadata in generated assets. & The current Python reference renderer is not a complete 3DTI substitute for publication-grade near-field/reverb rendering.\\
Consensus searches & Independently checked each assumption: burst/onset salience, HRTF/ILD/ITD primacy, reverb/DRR/IACC relevance, near-field importance, PPS tactile RT facilitation, individualized HRTF value, and the insufficiency of intensity alone. & The searches do not replace the need for local signal QC and behavioral validation; they justify the generation protocol, not the success of any future dataset.\\
\bottomrule
\end{longtable}

\section*{Consensus Search Export}

Manual Consensus searches were first run through the user's logged-in browser session on 2026-06-29. After the Consensus MCP bridge was verified locally, the same assumptions were rechecked through the read-only \texttt{Consensus Academic Research v3.4.0} MCP search tool. The searches were designed to test the specific assumptions behind adopting DynaSpace-style burst-train looming while keeping PPS-kit's 3DTI/SOFA renderer responsible for binaural spatial cues. Browser search pages and MCP-returned paper URLs are listed so the search trail can be re-opened later.

\begin{longtable}{P{0.21\linewidth} P{0.31\linewidth} P{0.23\linewidth} P{0.17\linewidth}}
\caption{Consensus searches and implementation relevance.}\\
\toprule
Question tested & Consensus result & Critical caveat & PPS-kit decision\\
\midrule
\endfirsthead
\toprule
Question tested & Consensus result & Critical caveat & PPS-kit decision\\
\midrule
\endhead
Do approaching looming sounds with rising intensity facilitate faster detection or orienting responses compared with receding or constant sounds?
\newline\url{https://consensus.app/search/approaching-looming-sound-detection/FUKQWKFNSbicvSXmbMchkA/}
& The result was a moderate-to-strong Yes: approaching/rising-intensity sounds are usually more salient and often produce faster orienting or detection responses.
& Some multisensory or visual-dominant tasks show no special looming advantage once task structure and cue type are controlled.
& Keep a strong approach/rising-salience profile, but treat it as a salience-plus-distance manipulation unless matched-level controls are present.\\
\addlinespace
Are repeated broadband noise bursts or onset transients useful for auditory spatial localization compared with continuous narrowband sounds?
\newline\url{https://consensus.app/search/auditory-spatial-localization-cues/85JQWFkfR3-f9kbKCELhBA/}
& The result was Yes: repeated broadband bursts and onset transients usually improve localization relative to continuous narrowband stimuli because they expose richer spectral, binaural, and onset cues.
& Repetition is not always superior to a longer single broadband burst; benefit depends on duration, reverberation, and task.
& Adopt the DynaSpace/Hobeika burst-train source profile as the dry source for the smartphone replication profile.\\
\addlinespace
Are HRTF spectral cues, interaural level differences, and interaural time differences the main cues for binaural spatial hearing and localization?
\newline\url{https://consensus.app/search/binaural-spatial-hearing-cues/3cA7IcmZQD2Vu5c-vJDfLw/}
& Consensus meter: 7 of 7 Yes. The evidence base treats HRTF spectral filtering, ILD, and ITD as the core binaural localization cues, with motion and reverberation adding support.
& Generic HRTFs can be adequate for horizontal lateralization but are weaker for elevation, front/back resolution, and individualized externalization.
& Do not imitate these cues by static spectral hacks in the dry WAV; render them through 3DTI/SOFA and record the chosen HRTF.\\
\addlinespace
Do reverberation, interaural coherence, and direct-to-reverberant ratio affect auditory externalization and distance perception in spatial audio?
\newline\url{https://consensus.app/search/auditory-externalization-spatial-audio/ZupUZceiRGGFHlFGafdDDQ/}
& Consensus meter: 10 of 10 Yes. Reverberation, IACC/coherence structure, and DRR all affect externalization and auditory distance.
& Diotic or poorly structured reverberation can support distance without producing convincing externalization; reverb must be spatially meaningful.
& Add reverb/DRR/IACC as explicit optional QC targets, not as untracked stereo decorrelation.\\
\addlinespace
Are near-field binaural cues important for auditory distance perception and localization for sources within one meter of the head?
\newline\url{https://consensus.app/search/near-field-binaural-auditory-cues/Fz9jdMNQS3e8cbWB0ziT3Q/}
& Consensus meter: 16 papers, 88 percent Yes and 13 percent Mixed. Near-field cues matter under 1 m, especially lateral ILD changes, but are not sufficient alone.
& The strongest near-field effects are lateral/off-midline and interact with level and DRR; ITD changes little with distance.
& Keep the 20 cm endpoint, but prefer native 3DTI near-field-capable rendering for publication stimuli.\\
\addlinespace
Do dynamic approaching sounds facilitate tactile reaction time in peripersonal-space tasks compared with receding or fixed sounds?
\newline\url{https://consensus.app/search/dynamic-sounds-and-tactile-reaction/dw8vS7TcQheii5_C_uOoNQ/}
& The result was Yes with caveats: approaching sounds usually speed tactile responses more than receding or fixed sounds in PPS tasks; prior summaries estimate near-far benefits on the order of tens of milliseconds.
& Expectancy and timing structure can inflate apparent PPS boundaries; Hobeika/Holmes-style critiques require clean controls.
& Use the DynaSpace looming-vs-fixed contrast for smartphone replication, but include matched-level and expectancy-aware variants for stronger inference.\\
\addlinespace
Does individualized or high-quality HRTF rendering improve spatial audio localization and externalization compared with generic non-individual HRTFs?
\newline\url{https://consensus.app/search/individualized-hrtf-spatial-audio/2eIYeqrxTeOiKpEOs8tCNg/}
& The result was mostly Yes: individualized HRTFs improve localization and usually externalization, strongest for elevation, front/back, and VR-like spatial presence.
& Carefully selected generic HRTFs can approach individual performance for some horizontal tasks.
& FABIAN remains an acceptable redistributable default, but the profile should make HRTF identity visible and allow higher-quality/individualized HRTFs later.\\
\addlinespace
Is rising intensity alone sufficient to create auditory looming perception, or do spectral, binaural, and reverberation cues add important spatial information?
\newline\url{https://consensus.app/search/auditory-looming-perception-cues/kEdiFP76QhasleThXrXE-Q/}
& The result was No as a complete account: rising intensity can create classic looming/warning effects, but spectral, binaural, and reverberation cues add spatial information.
& Intensity is often the dominant cue, so a strong gain ramp can masquerade as spatial looming if controls are weak.
& Use rising intensity as one declared component of approach, while preserving independent 3DTI-rendered HRTF, ILD, ITD, near-field, and room cues.\\
\bottomrule
\end{longtable}

\subsection*{Consensus MCP Paper Export}

The following compact export records the read-only MCP searches. The bracket numbers are search-local paper numbers returned by Consensus for that query, not global bibliography numbers. Full abstracts were not copied into this public report; the export preserves query text, result counts, representative titles, and direct Consensus paper URLs.

\begin{longtable}{P{0.20\linewidth} P{0.24\linewidth} P{0.48\linewidth}}
\caption{Consensus MCP query export.}\\
\toprule
Evidence target & Query and count & Representative MCP-returned papers\\
\midrule
\endfirsthead
\toprule
Evidence target & Query and count & Representative MCP-returned papers\\
\midrule
\endhead
Looming salience and warnings
& \texttt{approaching looming sounds rising intensity reaction time detection orienting receding constant sounds}; 20 papers.
& [3] Bach et al. (2009), ``Looming sounds as warning signals: the function of motion cues,'' \url{https://consensus.app/papers/details/45a18fb1f74b5912b30c9465fbb620b3/?utm_source=claude_desktop}. [6] Neuhoff (2001), ``An Adaptive Bias in the Perception of Looming Auditory Motion,'' \url{https://consensus.app/papers/details/0e0c538196b35477ab39922740f634dd/?utm_source=claude_desktop}. [8] Ignatiadis et al. (2024), ``Threat-related corticocortical connectivity elicited by rapid auditory looms,'' \url{https://consensus.app/papers/details/2114c61c357f5e8c854bbc197c06b63a/?utm_source=claude_desktop}.\\
\addlinespace
Broadband bursts and localization
& \texttt{auditory localization broadband noise bursts onset transients narrowband continuous sounds}; 20 papers.
& [1] Middlebrooks (1992), ``Narrow-band sound localization related to external ear acoustics,'' \url{https://consensus.app/papers/details/cef5e85bb5ae529d9543673adde4bf82/?utm_source=claude_desktop}. [2] Kim et al. (2024), ``Median Plane Localization of Short Noise Bursts in the Presence of Horizontally Spread White Noise Masker,'' \url{https://consensus.app/papers/details/dade8599ab9b5fdb849828fd8f9de951/?utm_source=claude_desktop}. [4] Macpherson et al. (2000), ``Localization of brief sounds: effects of level and background noise,'' \url{https://consensus.app/papers/details/b3bdb5d148d550bd8851c3c821390133/?utm_source=claude_desktop}.\\
\addlinespace
Binaural/HRTF cue ownership
& \texttt{binaural spatial hearing HRTF interaural level difference interaural time difference sound localization}; 19 papers.
& [4] Carlini et al. (2024), ``Auditory localization: a comprehensive practical review,'' \url{https://consensus.app/papers/details/ce0b49c63a025ebaa3abe91f3c12f873/?utm_source=claude_desktop}. [3] Gutierrez-Parera et al. (2022), ``Interaural time difference individualization in HRTF by scaling through anthropometric parameters,'' \url{https://consensus.app/papers/details/f79647e2cfc358ef8f3d52a1350d9051/?utm_source=claude_desktop}. [1] Lee et al. (2025), ``Binaural Sound Event Localization and Detection based on HRTF Cues for Humanoid Robots,'' \url{https://consensus.app/papers/details/11d9437d97a45d93b74195c9a49b37c0/?utm_source=claude_desktop}.\\
\addlinespace
Reverb, IACC, DRR, externalization
& \texttt{reverberation interaural coherence direct-to-reverberant ratio auditory externalization distance perception}; 20 papers.
& [1] Choe et al. (2023), ``Impact of direct-to-reverberation ratio and interaural cross-correlation on externalization in binaural audio rendering,'' \url{https://consensus.app/papers/details/e1c815e6c9a55a8db398292d85a3c401/?utm_source=claude_desktop}. [2] Lavandier et al. (2025), ``Reverberation cues underlying virtual auditory distance perception for a frontal source,'' \url{https://consensus.app/papers/details/dd37a407e06a51cbab6251c5ac60e131/?utm_source=claude_desktop}. [5] Kolarik et al. (2013), ``Discrimination of virtual auditory distance using level and direct-to-reverberant ratio cues,'' \url{https://consensus.app/papers/details/9b0b8eedb09d5eff8fbcca667b091eec/?utm_source=claude_desktop}.\\
\addlinespace
Near-field distance cues
& \texttt{near-field binaural cues auditory distance perception within one meter ILD ITD}; 20 papers.
& [2] Brungart et al. (1999), ``Auditory localization of nearby sources. Head-related transfer functions,'' \url{https://consensus.app/papers/details/830a7bfa7d5c5bb4b88e4822debf326d/?utm_source=claude_desktop}. [4] Zhu et al. (2023), ``The contributions of level and near-field head-related transfer functions cues on auditory distance perception in a free field,'' \url{https://consensus.app/papers/details/cd5157c315985b23ab71320b317f770a/?utm_source=claude_desktop}. [5] Guo et al. (2023), ``Distance discrimination thresholds of proximal sound sources in a real anechoic environment,'' \url{https://consensus.app/papers/details/e2c5ce9ed8fd5672aec2443b05c4b141/?utm_source=claude_desktop}.\\
\addlinespace
PPS tactile RT phenomenon
& \texttt{peripersonal space dynamic approaching sounds tactile reaction time fixed receding}; 19 papers.
& [1] Canzoneri et al. (2012), ``Dynamic Sounds Capture the Boundaries of Peripersonal Space Representation in Humans,'' \url{https://consensus.app/papers/details/a947a34ca17b5728b5ae5429d47f8579/?utm_source=claude_desktop}. [2] Matsuda et al. (2021), ``Peripersonal space in the front, rear, left and right directions for audio-tactile multisensory integration,'' \url{https://consensus.app/papers/details/c34b54676bdf50eda323e11dc189b316/?utm_source=claude_desktop}. [3] Hobeika et al. (2019), ``Capturing the dynamics of peripersonal space by integrating sound propagation properties and expectancy effects,'' \url{https://consensus.app/papers/details/53d4e77477fd5a2cad2abd9ea1789dca/?utm_source=claude_desktop}.\\
\addlinespace
Individualized/high-quality HRTF value
& \texttt{individualized HRTF generic HRTF localization externalization spatial audio}; 20 papers.
& [1] Jenny et al. (2020), ``Usability of Individualized Head-Related Transfer Functions in Virtual Reality,'' \url{https://consensus.app/papers/details/f3acff26460c5eadba2bc04664a8a177/?utm_source=claude_desktop}. [2] Frankenbach et al. (2025), ``Evaluating Localization Accuracy of Dynamic Binaural Synthesis,'' \url{https://consensus.app/papers/details/ed8f8111a884542faaac73d52f5f30fa/?utm_source=claude_desktop}. [4] SADIE II case study (2018), ``A Perceptual Evaluation of Individual and Non-Individual HRTFs,'' \url{https://consensus.app/papers/details/c390fc1e332f5147a003b318974b735a/?utm_source=claude_desktop}.\\
\addlinespace
Intensity is not the whole spatial cue set
& \texttt{auditory looming rising intensity spectral cues binaural reverberation distance perception}; 20 papers.
& [3] Kolarik et al. (2015), ``Auditory distance perception in humans: a review of cues,'' \url{https://consensus.app/papers/details/fde5d69ba97851afac0015ca1cf942c9/?utm_source=claude_desktop}. [4] Baumgartner et al. (2017), ``Asymmetries in behavioral and neural responses to spectral cues demonstrate the generality of auditory looming bias,'' \url{https://consensus.app/papers/details/40f6ff7efbcc50f980772b493696877b/?utm_source=claude_desktop}. [5] Mershon et al. (1975), ``Intensity and reverberation as factors in the auditory perception of egocentric distance,'' \url{https://consensus.app/papers/details/af28b552a18a54db9113a715842b4601/?utm_source=claude_desktop}.\\
\bottomrule
\end{longtable}

\subsection*{Consensus-Derived Implementation Claims}

\begin{enumerate}
  \item The dry source should be a burst train. This is supported by the raw DynaSpace files, the Hobeika lineage, and Consensus evidence that broadband transients support localization and attentional salience.
  \item Spatialization should be delegated to 3DTI/SOFA rather than baked into a static waveform. Consensus strongly supports HRTF, ILD, ITD, near-field, reverberation, and DRR as spatial cues that belong in the renderer.
  \item The original smartphone fixed-vs-looming contrast remains a replication target, but it is not a clean causal design by itself. Rising intensity and fixed-source level differences can drive response facilitation unless matched-level control variants are included.
  \item Reverb and low interaural coherence are plausible missing externalization features in the current proxy, but they should be implemented as explicit room/BRIR parameters with QC output.
  \item The toolkit should preserve trajectory flexibility: DynaSpace's 45-degree left approach is one profile, not a hard-coded spatial renderer.
\end{enumerate}


\section*{Measured Raw Signal Inventory}

Measurements were generated by \path{measure_dynaspace_audio.py}. Burst detection used a 20 ms RMS envelope, 2.5 ms hop, 3 dB peak prominence, and at least 55 ms spacing. Binaural lag was measured by cross-correlation within +/-1 ms; the sign is a script convention and should not be treated as a biological left/right truth until a channel-sign probe is run.

\begin{table}[h]
\centering
\small
\caption{Compact derived feature summary.}
\begin{tabular}{lrrrr}
\toprule
Metric & Dyna loom & Dyna fixed & PPS loom & PPS fixed\\
\midrule
Sample rate (Hz) & 48000 & 48000 & 44100 & 44100\\
Duration (s) & 3.850 & 3.850 & 3.850 & 3.850\\
Peak (dBFS) & -3.19 & -33.48 & -0.91 & -0.91\\
RMS (dBFS) & -37.28 & -53.25 & -31.33 & -24.64\\
Crest factor (dB) & 34.10 & 19.77 & 30.42 & 23.73\\
Detected burst peaks & 33 & 33 & 33 & 33\\
Median burst IOI (ms) & 95.0 & 95.0 & 94.8 & 94.8\\
Active p90-p10 level span (dB) & 18.20 & 8.09 & 33.66 & 42.20\\
Active level slope (dB/s) & 4.31 & -1.16 & 8.33 & -0.15\\
Spectral centroid (Hz) & 10948 & 9319 & 8107 & 8173\\
95\% rolloff (Hz) & 19910 & 19846 & 15515 & 15601\\
ILD L-R (dB) & -6.34 & -4.26 & 12.61 & 12.70\\
Cross-corr lag (us) & 187.5 & 187.5 & -362.8 & -362.8\\
IACC abs & 0.229 & 0.132 & 0.703 & 0.690\\
Moving 100 ms IACC range & 0.043--0.269 & 0.043--0.228 & 0.647--0.751 & 0.646--0.727\\
\bottomrule
\end{tabular}
\end{table}

\begin{figure}[h]
\centering
\includegraphics[width=0.92\linewidth]{figures/rms_envelope_comparison.png}
\caption{20 ms RMS envelope comparison. DynaSpace raw looming is pulsed and rises over the trial; the regenerated PPS-kit proxy now uses the same burst-train timing but still carries a stronger level ramp.}
\end{figure}

\begin{figure}[h]
\centering
\includegraphics[width=0.92\linewidth]{figures/spectral_band_comparison.png}
\caption{Spectral band comparison. The DynaSpace raw looming file is brighter, with more relative high-frequency energy and a higher centroid than the current PPS-kit proxy.}
\end{figure}

\begin{figure}[h]
\centering
\includegraphics[width=0.72\linewidth]{figures/binaural_summary.png}
\caption{Binaural summary. The raw DynaSpace looming file has lower interaural coherence and opposite-signed/much smaller ILD than the current PPS-kit proxy.}
\end{figure}

\clearpage
\section*{One-to-One Feature Audit}

\begin{longtable}{P{0.15\linewidth} P{0.20\linewidth} P{0.23\linewidth} P{0.22\linewidth} P{0.13\linewidth}}
\caption{Feature-by-feature decision matrix.}\\
\toprule
Feature & Raw DynaSpace evidence & Spatialization or PPS rationale & PPS-kit / 3DTI status & Decision\\
\midrule
\endfirsthead
\toprule
Feature & Raw DynaSpace evidence & Spatialization or PPS rationale & PPS-kit / 3DTI status & Decision\\
\midrule
\endhead
Broadband Gaussian/noise source & Raw files are broadband stereo noise bursts. DynaSpace looming centroid is 10.95 kHz and 95\% rolloff is 19.91 kHz. & Broadband content makes HRTF spectral filtering, ILD, and localization cues more available than narrowband tones \cite{carlini2024}. It also matches the Hobeika PPS lineage \cite{hobeika2020}. & The regenerated PPS proxy now uses generated white-noise bursts, but its spectrum remains less bright than the raw DynaSpace file. 3DTI can spatialize any dry source; the dry source spectrum should be controlled before rendering. & \statusmodify\\
Burst train temporal structure & Raw looming and fixed files both contain 33 detected peaks, median IOI 95 ms. Measured half-prominence widths are 55 ms looming and 70 ms fixed, inflated by ramps/reverb/detector relative to dry burst design. & The burst train creates repeated onsets while preserving a continuous trial. Onsets are strong for alerting and localization, and the 95 ms spacing agrees with the Hobeika lineage. & The regenerated PPS proxy now has 33 detected bursts and about 94.8 ms median IOI. This feature is adopted as the default source profile. & \statuskeep\\
Burst ramps and click control & The lineage uses short burst ramps; measured raw files have high crest factors but no clipping. & Ramps avoid click artifacts while retaining salient attacks. & PPS generator should create dry bursts with explicit raised-cosine rise/fall parameters, not rely on random noise fluctuations. & \statuskeep\\
Overall duration and silence window & Raw files are 3.85 s. Paper reports a 3250 ms auditory stimulus preceded by 300 ms silence and followed by a randomized 1000--1400 ms silence at trial level \cite{roussel2025}. & Timing window matters because tactile anchors sample distance over the approach. & PPS profile matches 3.85 s proxy duration and anchor logic. Keep metadata explicit because paper-level and raw-file-level timings are not identical descriptions. & \statusmodify\\
Distance trajectory & Paper trajectory: 640 cm to 20 cm at ear level, constant speed 210 cm/s, left hemispace azimuth 45 degrees; anchors map to 640, 320, 160, 80, 40, and 20 cm \cite{roussel2025}. & This is the main PPS geometry. It separates temporal expectancy from spatial proximity through fixed-vs-looming contrast. & PPS template already encodes 6.4 m to 0.2 m, 2.105 m/s, signed -45 degrees, and anchors. 3DTI is better than a baked file for future angles. & \statusalready\\
Rising intensity / distance gain & DynaSpace raw looming active p90-p10 span is 18.20 dB; regenerated PPS proxy is 33.66 dB. Raw fixed is much quieter overall than raw looming. & Rising level is a core auditory looming warning cue \cite{bach2008,neuhoff2016}; distance intensity is also a primary distance cue \cite{kolarik2016,carlini2024}. & PPS proxy still has a stronger gain ramp and higher RMS than the raw looming file. Current Python path applies 1/r amplitude gain when gain law is enabled, then peak-normalizes unless loudness policy disables it. Native 3DTI can express direct-path attenuation more cleanly. & \statusmodify\\
Fixed comparator level & Raw fixed peak is -33.48 dBFS and RMS is -53.25 dBFS, far lower than raw looming. & A far fixed comparator may intentionally represent source distance, but it also risks making the main effect partly a loudness/salience contrast. & PPS fixed proxy now shares burst timing but remains much louder than the raw fixed file. For exact replication, preserve far-level comparator; for clean inference, add a matched-level control profile. & \statusskip\\
High-frequency brightness & DynaSpace looming centroid is 10.95 kHz; regenerated PPS proxy is 8.11 kHz. Raw 95\% rolloff is near 19.9 kHz, proxy near 15.5 kHz. & High frequencies strengthen ILD and pinna/HRTF spectral cues; broad bandwidth improves localization \cite{carlini2024}. & 3DTI should supply HRTF coloration, but the dry source should remain bright enough. Add optional spectrum matching/shelf only as a dry-source setting, not as a fake spatializer. & \statusmodify\\
HRTF spectral filtering & Paper reports Max/MSP Spat and non-individual LISTEN HRTFs for direct sound and early reflections \cite{roussel2025}. & HRTFs carry head/pinna/torso spectral cues for externalized spatial location \cite{carlini2024}. & PPS-kit uses FABIAN SOFA as a redistributable standard; LISTEN is excluded pending license audit. Full 3DTI/SOFA rendering is a stronger, parameterized approach than baking the raw spectral coloration. & \statusalready\\
ILD & DynaSpace looming global ILD L-R is -6.34 dB; moving-window ILD ranges -7.78 to 0.48 dB. PPS proxy global ILD is +12.61 dB and nearly static. & ILD is a main lateralization cue, especially at high frequencies and near-field positions \cite{carlini2024}. & 3DTI should implement ILD through HRTF, distance, head shadow, and near-field corrections. Current proxy sign/magnitude mismatch means coordinate/channel convention must be audited before using the profile for left/right claims. & \statusmodify\\
ITD / cross-ear lag & DynaSpace raw lag is +187.5 us; PPS proxy lag is -362.8 us by the same script convention. & ITD is a principal horizontal cue, especially for lower-frequency components; broadband bursts provide onset information too \cite{carlini2024}. & Full 3DTI supports customized ITD and propagation. The Python proxy uses nearest HRIR frames and does not prove correct signed ITD. Add a channel-sign validation file. & \statusmodify\\
Interaural coherence / IACC & DynaSpace looming IACC abs is 0.229 globally and 0.043--0.269 in 100 ms windows. PPS proxy is much more coherent at about 0.703 globally. & Low interaural coherence is consistent with reflected/reverberant binaural energy and can aid externalization, but excessive decorrelation can weaken precise lateral localization. & 3DTI with BRIR/Ambisonic reverb is a better controlled way to add this than arbitrary decorrelation. Current PPS profile has reverb disabled. & \statusmodify\\
Early reflections and late reverb & Paper reports direct sound, early reflections, and late reverberation in the smartphone stimuli \cite{roussel2025}. Raw low IACC supports the presence of non-anechoic spatial rendering. & Reverberation contributes distance/externalization information; direct-to-reverberant ratio is a major distance cue \cite{kolarik2016,carlini2024}. & Full 3DTI supports BRIR/Ambisonic reverberation \cite{cuevas2019}; current PPS Python reference config declares reverb disabled. Add an optional DynaSpace-room profile rather than always-on reverb. & \statuskeep\\
Direct-to-reverberant ratio & Not directly measured from the stereo WAV without the dry/render components, but low IACC and paper methods imply a baked direct-plus-reflection mixture. & DRR can provide distance information that is less dependent on source level than raw intensity \cite{kolarik2016}. & 3DTI can model this better if a room/BRIR configuration is explicit. Current proxy has no explicit DRR target. & \statusmodify\\
Near-field distance cues & Endpoint is 20 cm, inside near field. & Near-field ILD/head-shadow changes become important close to the head; source distance changes are perceptually larger near the body. & 3DTI full high-quality path includes near-field corrections \cite{cuevas2019}; current Python proxy is a simplified reference render. Use native 3DTI for publication-grade near-field stimuli. & \statusalready\\
Trajectory flexibility & Raw WAV is a single baked 45-degree left-hemispace stimulus. & Future studies need different approach angles and maybe curved trajectories. & PPS-kit/3DTI is clearly better here because the trajectory is metadata-driven. Do not encode spatial path as immutable audio. & \statusalready\\
Sample rate / bit depth & DynaSpace raw files are 48 kHz, 16-bit PCM. PPS proxy files are 44.1 kHz. Paper text reports 44.1 kHz source preparation, while Android raw files are 48 kHz. & Sample rate affects top-band content and exact asset playback, but it is not itself a spatial cue. & Match the HRTF/SOFA/render chain sample rate and document resampling. Use 48 kHz only for exact Android playback replication, not because it is inherently better. & \statusskip\\
Peak/RMS/crest quality & Raw looming peak -3.19 dBFS, no clipping, RMS -37.28 dBFS, crest 34.10 dB. Regenerated PPS proxy peak -0.91 dBFS, RMS -31.33 dBFS, crest 30.42 dB. & High crest is expected for sparse burst trains. Hidden peak normalization can destroy intended absolute or relative level policies. & PPS should disable hidden peak normalization in calibrated mode and store SPL/RMS targets. Raw DynaSpace quality is usable, but its fixed-vs-looming level asymmetry must be declared. & \statusmodify\\
Phone/transducer chain & Paper used an Android phone and headphones; the raw file alone does not encode the full playback chain. & Smartphone output and headphone response can affect level and spectrum, but this is apparatus, not a spatial rendering principle. & PPS-kit should not mimic unknown phone coloration unless validating Android playback specifically. Use measured audio-interface/headphone calibration instead. & \statusskip\\
Tactile vibration timing & Paper vibration: 155 Hz, about 50 ms; source profile records Android code timing details. & This affects PPS RT alignment, not audio spatial quality. & Keep DynaSpace timing anchors and vibration metadata in the study profile, but do not mix haptic waveform design into audio spectral matching. & \statusmodify\\
\bottomrule
\end{longtable}

\section*{Signal Quality Assessment}

The DynaSpace raw looming WAV is technically clean for a behavioral stimulus: it has no clipping, large headroom, high crest factor consistent with a pulse train, and a stable 33-burst temporal pattern. Its level is low on average because most of the file is sparse/noisy and because distance/reverberation were baked into the render. The raw fixed file is also clean, but its peak and RMS are much lower than the looming file. That fixed-vs-looming level gap is not a file-quality problem; it is a design property that can become an interpretive problem if a future study tries to attribute all facilitation to spatial proximity rather than level/salience.

The PPS-kit proxy files are also technically clean. After this implementation pass they now match the core DynaSpace burst-train timing: 33 detected bursts and about 95 ms median IOI. They are still not fully DynaSpace-equivalent: they are peak-normalized close to full scale, the looming proxy has a stronger gain ramp than the raw DynaSpace looming file, the dry/source spectrum is less bright, and the proxy binaural image is much more coherent. This makes the present proxy a stronger smartphone-style source approximation, but not a bit-matched recreation of the original Spat/LISTEN/room/phone render.

\section*{Best-of-Both-Worlds Generation Protocol}

\begin{enumerate}
  \item Generate a dry, mono Gaussian white-noise burst train: 33 bursts, 30 ms dry burst duration, 10 ms raised-cosine rise/fall, 65 ms silent interval, median peak spacing 95 ms, seed recorded in metadata. Keep a source-spectrum option that can target the raw DynaSpace brightness without pretending that spectrum matching is spatial rendering.
  \item Use the DynaSpace trajectory profile already encoded in PPS-kit: 6.4 m to 0.2 m, 2.105 m/s, ear level, constant-bearing left approach, tactile anchors at 105, 1625, 2385, 2765, 2955, and 3050 ms. Preserve the ability to change azimuth/elevation/endpoints for future studies.
  \item Render with native 3DTI high-quality mode when available. Use SOFA HRIRs, HRIR interpolation, customized ITD/head model, propagation delay, 6 dB-per-distance-doubling direct-path attenuation, and near-field corrections. Use FABIAN as the redistributable default; use LISTEN only in private or after a license audit.
  \item Add an optional DynaSpace-room/externalization layer: early reflections plus late reverberation or BRIR/Ambisonic reverb. Tune it with QC targets rather than by ear: lower IACC than the dry proxy, documented direct-to-reverberant settings, and retained direct-source lateralization.
  \item Provide two comparator modes. A \textit{DynaSpace replication} mode keeps the far fixed source at 640 cm with the original-style distance/level logic. A \textit{clean causal} mode adds matched-level controls so spatial proximity, temporal expectancy, and loudness/salience can be separated.
  \item Add render QC acceptance checks for every generated WAV: no clipping; declared peak/RMS/SPL policy; 33 detected bursts and 95 ms median IOI; active looming level span near the chosen target; channel-sign validation; expected ILD/ITD direction for the requested azimuth; IACC/DRR report if reverb is enabled; and all trajectory anchors stored in sidecar JSON.
\end{enumerate}

\section*{Implementation Consequences for PPS-kit}

\begin{itemize}
  \item PPS-kit should treat the best-of-both-worlds protocol as the new default operating standard for ops-generated looming stimuli: a DynaSpace/Hobeika-style broadband burst-train source, explicit level policy, and 3DTI/SOFA-owned spatial rendering.
  \item New generated-noise looming bakes should default to a named \texttt{dynaspace\_gaussian\_burst\_train} source profile: 33 bursts, 30 ms burst duration, 10 ms rise/fall, 65 ms inter-burst interval, and 300 ms onset. Published-study profiles may override this only when the source paper specifies a different waveform or apparatus constraint.
  \item The DynaSpace profile should expose two level-policy presets: \texttt{smartphone\_replication} and \texttt{matched\_level\_control}.
  \item The renderer should prefer native 3DTI for DynaSpace-style publication stimuli; the Python SOFA reference path is acceptable for preview/proxy but should be flagged as simplified.
  \item Reverb/BRIR should be optional but first-class. The original smartphone stimulus used reflections/reverb, and the measured low IACC shows that current PPS proxy externalization is not equivalent.
  \item Before changing left/right signs, run a channel-orientation audit. The raw files and current proxy disagree in ILD/ITD sign, and that can be a coordinate convention issue rather than a true stimulus mismatch.
\end{itemize}

\section*{Golden-Standard Operating Protocol}

The adopted standard is not ``copy the phone WAV.'' It is source-renderer separation:

\begin{enumerate}
  \item \textbf{Dry source}: generate a salient, bright, broadband Gaussian white-noise burst train using the DynaSpace/Hobeika timing defaults unless the target study has a documented different source.
  \item \textbf{Trajectory}: store approach geometry as editable metadata: start/end distance, speed, azimuth/elevation, endpoint holds, and SOA/distance anchors.
  \item \textbf{Spatial rendering}: use 3DTI/SOFA for HRTF spectral filtering, ILD, ITD, distance attenuation, propagation delay, near-field behavior, and future trajectory angles.
  \item \textbf{Externalization}: add room/BRIR/reverb/DRR/IACC only through explicit renderer parameters and QC outputs, not hidden stereo decorrelation.
  \item \textbf{Controls}: provide both replication-like high-salience contrasts and matched-level causal controls wherever the research question needs to distinguish PPS proximity from raw salience.
  \item \textbf{QC}: every generated looming source should report no clipping, source-profile metadata, burst count/IOI, peak/RMS/level policy, trajectory anchors, channel-sign/azimuth expectations, and room metrics when reverb is enabled.
\end{enumerate}

\section*{Bottom Line}

It is worth incorporating several DynaSpace stimulus features into PPS-kit, but only as explicit, testable generation parameters. The strongest additions are the burst-train source, DynaSpace-level trajectory anchors, controlled rising level, and optional reflection/reverb/externalization layer. The features that 3DTI should own are spatial cues: HRTF spectrum, ILD, ITD, near-field behavior, propagation delay, distance attenuation, and arbitrary future trajectories. The raw smartphone WAV is best treated as an audit target and validation reference, not as the renderer.

\clearpage
\begin{thebibliography}{99}

\bibitem{roussel2025}
U. Roussel, E. Flety, C. Agon, I. Viaud-Delmon, and M. Taffou.
\textit{Using Android Smartphones to Collect Precise Measures of Reaction Times to Multisensory Stimuli}.
Sensors 25(19):6072, 2025.
\url{https://doi.org/10.3390/s25196072}

\bibitem{dynaspace_repo}
DynaSpace source repository.
\url{https://forge-2.ircam.fr/roussel/dynaspace/}

\bibitem{hobeika2020}
L. Hobeika, M. Taffou, I. Viaud-Delmon, and colleagues.
\textit{Capturing the dynamics of peripersonal space by integrating expectancy effects and sound propagation properties}.
Journal of Neuroscience Methods, 2020.
\url{https://doi.org/10.1016/j.jneumeth.2019.108534}

\bibitem{canzoneri2012}
E. Canzoneri, G. Magosso, and A. Serino.
\textit{Dynamic Sounds Capture the Boundaries of Peripersonal Space Representation in Humans}.
PLOS ONE 7(9):e44306, 2012.
\url{https://doi.org/10.1371/journal.pone.0044306}

\bibitem{bach2008}
D. R. Bach, H. Schachinger, J. G. Neuhoff, F. Esposito, F. Di Salle, C. Lehmann, M. Herdener, K. Scheffler, and E. Seifritz.
\textit{Rising sound intensity: an intrinsic warning cue activating the amygdala}.
Cerebral Cortex 18(1):145--150, 2008.
\url{https://doi.org/10.1093/cercor/bhm040}

\bibitem{neuhoff2016}
J. G. Neuhoff and colleagues.
\textit{Looming sounds are perceived as faster than receding sounds}.
Cognitive Research: Principles and Implications, 2016.
\url{https://doi.org/10.1186/s41235-016-0017-4}

\bibitem{carlini2024}
A. C. Carlini, C. Bordeau, M. Ambard, and colleagues.
\textit{Auditory localization: a comprehensive practical review}.
Frontiers in Psychology, 2024.
\url{https://doi.org/10.3389/fpsyg.2024.1408073}

\bibitem{kolarik2016}
A. J. Kolarik, B. C. J. Moore, P. Zahorik, S. Cirstea, and S. Pardhan.
\textit{Auditory distance perception in humans: a review of cues, development, neuronal bases, and effects of sensory loss}.
Attention, Perception, and Psychophysics, 2016.
\url{https://pmc.ncbi.nlm.nih.gov/articles/PMC4744263/}

\bibitem{cuevas2019}
M. Cuevas-Rodriguez, L. Picinali, A. Gonzalez-Toledo, C. Garre, E. de la Rubia-Cuestas, L. Molina-Tanco, and A. Reyes-Lecuona.
\textit{3D Tune-In Toolkit: An open-source library for real-time binaural spatialisation}.
PLOS ONE 14(3):e0211899, 2019.
\url{https://doi.org/10.1371/journal.pone.0211899}

\bibitem{sofa}
SOFA Spatially Oriented Format for Acoustics project and database.
\url{https://www.sofaconventions.org/}

\bibitem{fabian}
FABIAN neutral HRIR resource, TU Berlin / DepositOnce.
\url{https://depositonce.tu-berlin.de/items/3b423df7-a764-4ce1-9065-4e6034bba759}

\end{thebibliography}

\end{document}
